We first prove the upper bound.  Fix an arm \(t\), a boundary point \(x\), and a bandwidth \(0<h\leq r_0\).  Write \(R_x=\lVert X-x\rVert_2\) and
\[
 W_{t,x,h}=h^{-2}\mathbf1\{T=t\}K_\triangle(R_x/h).
\]
For each fixed \(x\), absolute continuity gives \(P(X=x)=0\), so replacing the strict-sign summand in the estimator by \(W_{t,x,h}\) does not change its population moments.  The stability lemma gives
\[
 d_{t,P,x}(h):=\mathbb E_PW_{t,x,h}\geq\kappa_0.
\]
The upper-mass lemma, \(0\leq K_\triangle\leq1\), and the bounded regression imply
\[
 \mathbb E_PW_{t,x,h}^2\leq \pi Lh^{-2},
 \qquad
 \mathbb E_P\{W_{t,x,h}Y\}^2\leq Ch^{-2}.
\]
For the second inequality, condition on \(X\), use consistency on arm \(t\), the uniform bound on \(\mu_{t,P}\), and the conditional sub-Gaussian second-moment bound.

The same conditioning gives, for every integer \(k\geq2\),
\[
 \mathbb E_P|W_{t,x,h}\{Y-\mu_{t,P}(X)\}|^k
 \leq k!C^k h^{2-2k}.
\]
Expanding the exponential series and applying exponential Markov inequality therefore yields constants \(c,C>0\), uniform in \(P,t,x\), such that for either the centered denominator or centered numerator average,
\[
 P\!\left(\left|(\mathbb P_n-P)Z_{t,x,h}\right|>u\right)
 \leq2\exp\!\left\{-\frac{cnu^2}{h^{-2}+h^{-2}u}\right\}.
\]
This is the required Bernstein inequality written with its variance and scale terms; no limit law is assumed.

Cover \(\mathcal B_P\subseteq[-L,L]^2\) by an ambient Euclidean \(\epsilon_n\)-net with \(\epsilon_n=n^{-4}\).  Its cardinality is at most \(C_Ln^8\), independently of the geometry of \(\mathcal B_P\).  The map \(x\mapsto K_\triangle(\lVert X_i-x\rVert_2/h)\) is \(h^{-1}\)-Lipschitz.  Thus the normalized denominator and numerator oscillations between a point and its nearest net point are at most
\[
 C\epsilon_nh^{-3},
 \qquad
 C\epsilon_nh^{-3}\mathbb P_n|Y|,
\]
respectively, and the corresponding population oscillations obey the same bounds with \(\mathbb E_P|Y|\) in place of \(\mathbb P_n|Y|\).  The actual estimator excludes an observation when \(X_i=x\).  Since an absolutely continuous sample has no ties almost surely, this changes the normalized denominator by at most \((nh^2)^{-1}\) and the normalized numerator by at most \((nh^2)^{-1}\max_i|Y_i|\), uniformly in \(x\).  Conditional sub-Gaussianity and the bounded regression give
\[
 \sup_P\mathbb E_P\max_{i\leq n}|Y_i|^2\leq C\log n,
 \qquad
 \sup_P\mathbb E_P(\mathbb P_n|Y|)^2\leq C.
\]
Consequently, a union bound over the net, followed by the displayed oscillation bounds, shows that for every fixed \(A_0>0\), after enlarging \(C\), with probability at least \(1-Cn^{-A_0}\),
\[
 \sup_{t,x}\left|\mathbb P_nW_{t,x,h}-d_{t,P,x}(h)\right|
 +\sup_{t,x}\left|\mathbb P_n(W_{t,x,h}Y)-\mathbb E_P(W_{t,x,h}Y)\right|
 \leq C\left\{\sqrt{\frac{\log n}{nh^2}}+\frac{\log n}{nh^2}\right\},
\]
provided \(h\asymp a_n\).  All suprema here range over the two arms and \(x\in\mathcal B_P\).

For large \(n\), the right side is smaller than \(\kappa_0/2\), so every empirical denominator is positive and bounded below by \(\kappa_0/2\).  The population ratio is a kernel-weighted average of \(\mu_{t,P}(z)\) over \(R_x\leq h\).  Since \(q\geq1\),
\[
 \left|\frac{\mathbb E_P(W_{t,x,h}Y)}{d_{t,P,x}(h)}-\mu_{t,P}(x)\right|
 \leq Lh.
\]
The elementary ratio identity
\[
 \frac{\widehat r}{\widehat d}-\frac r d
 =\frac{\widehat r-r}{\widehat d}
 +\frac r d\frac{d-\widehat d}{\widehat d}
\]
therefore gives, on the good event,
\[
 \sup_{x\in\mathcal B_P}
 |\widehat\tau^{\mathrm{LC}}_{n,h}(x)-\tau_P(x)|
 \leq C\left\{h+\sqrt{\frac{\log n}{nh^2}}+\frac{\log n}{nh^2}\right\}.
\]
The causal target is not defined by this ratio: the canonical exact-distance lemma separately identifies its arm limits and the preceding display bounds the estimator of those limits.

On the exceptional event each arm estimate is either zero or a convex combination of observed outcomes, so the supremum loss is at most \(2\max_i|Y_i|+2L\).  Cauchy--Schwarz and the preceding maximal second-moment bound make its expected contribution negligible when \(A_0\) is chosen large.  With \(h\asymp a_n\), all three terms in braces are bounded by \(Ca_n\).  The estimator is a measurable member of \(\mathcal T_n^{\mathrm{dist}}\) by its explicit construction.  Hence
\[
 R_n^{\mathrm{dist}}(\mathcal P_{\mathrm{rect}})\leq C_{mathrm{rect}}a_n.
\]

For the converse, use the hypercube in the angular-density packing.  Poissonize the sample at mean \(2n\), attach independent continuous marks, and partition the Poisson process into the disjoint bump cells and their complement.  Conditional on the hypercube bits, these restrictions are independent, the complement is bit-independent, and the law in cell \(j\) depends only on bit \(j\).  Decoder \(j\), obtained by midpoint-thresholding \(\widehat\tau_n(x_j)\), observes its local cell only through the compressed statistic \(U_{n,P}(x_j)\), together with the other cells and the bit-independent complement.  Poisson tensorization multiplies the one-observation divergence by \(2n\); the slack \(\alpha<1/8\) in the packing makes the resulting bound \(2\alpha\log M_n\) with \(2\alpha<1\).  The coordinatewise-overlap direct-product lemma thus gives
\[
 \Pr\{\text{at least one bit is decoded incorrectly}\}\geq\frac12-o(1).
\]
If the boundary sup-norm loss were smaller than \(c_1a_n/2\), every midpoint decoder would be correct.  Hence the Poissonized expected loss is at least \(c a_n\).  Given a Poisson sample of mean \(2n\), marked subsampling produces exactly \(n\) independent observations on the event that its size is at least \(n\); the complement has probability at most \(e^{-cn}\).  Thus any estimator for exactly \(n\) observations would induce a Poissonized estimator with the same risk up to a negligible term.  It follows that
\[
 R_n^{\mathrm{dist}}(\mathcal P_{\mathrm{rect}})\geq c_{mathrm{rect}}a_n.
\]
Dividing the two inequalities by \(a_n\) and taking \(\liminf\) and \(\limsup\) proves the theorem.  The construction belongs to the project-defined class proved above; no inclusion in the class of the cited comparator theorem is used or asserted.